\section{Theoretical setup}

In this section we briefly introduce the nucleon matrix elements for the quasi-PDFs that we aim to renormalize.
We also explain the renormalization prescription for the three types of PDFs: unpolarized, helicity and transversity.
For details on the computation of the nucleon matrix elements, we refer to Refs.~\cite{Alexandrou:2015rja,Alexandrou:2016jqi}.

%#############################################################
\vspace*{0.25cm}
\subsection{Nucleon matrix elements}
\vspace*{0.15cm}
%#############################################################

We consider matrix elements of non-local fermion operators that contain a straight Wilson line, denoted by
$h_\Gamma(P_3,z)$. The variable $z$ is the length of the Wilson line, and $P_3$ is the nucleon momentum,
which is taken in the same direction as the Wilson line.
The quasi-PDFs can be computed from the Fourier transform of the following local matrix elements:
\begin{equation}
\label{MElements}
h_\Gamma(P_3,z)=\langle N |\bar{\psi}(0,z)\,\Gamma \,W_3(z) \psi(0,0) |N\rangle,
\end{equation}
where $|N\rangle$ is a nucleon state with spatial momentum $\vec{P}{=}(0,0,P_3)$ along the $3$-direction and
$W_3(z)$ is a Wilson line of length $z$ in the same direction. $\Gamma$ denotes the Dirac structure of the
operator insertion, which is $\gamma_\mu$ (unpolarized), $\gamma_\mu\cdot\gamma_5$ (helicity), $\sigma_{\mu \nu}$
(transversity). In the works appearing in the literature $\gamma_\mu$ is taken along the Wilson line. In principle, one may
choose $\gamma_\mu$ orthogonal to the direction of the Wilson line.  In this case, the unpolarized operator is
free of mixing, while the helicity and transverity do mix. For example, choosing the $\gamma$-matrix in the temporal direction
is important for a faster convergence to the physical PDFs, as discussed in Ref.~\cite{Radyushkin:2016hsy}. Our recent
work~\cite{Constantinou:2017sej} indicates that $\vec P_3 {\perp} \vec z$ \,\,($\vec P_3 {\parallel} \vec z$) is ideal for
the unpolarized (helicity and transversity) case.

\medskip
To calculate the bare matrix elements, we use the setup of Ref.~\cite{Alexandrou:2016jqi}.
We consider one ensemble of dynamical $N_f{=}2{+}1{+}1$ twisted mass fermions produced by ETMC~\cite{Baron:2010bv}, with volume
$32^3 {\times} 64$, lattice spacing $a {\approx}0.082$\,fm~\cite{Carrasco:2014cwa} and a bare twisted mass of $a\mu {=} 0.0055$,
which corresponds to a pion mass of around 375 MeV. We performed our calculations on 1000 gauge configurations with 15
forward propagators and 2 stochastic propagators, i.e.\ 30000 measurements in total. We will present results for momentum $P_3=\frac{6\pi}{L}$,
which is around 1.4 GeV in physical units. Gaussian smearing has been employed on the nucleon interpolating fields in the calculation of the matrix elements
\cite{Alexandrou:2016jqi}.

%#############################################################
\vspace*{0.25cm}
\subsection{Renormalization scheme}
\label{sec:scheme}
\vspace*{0.15cm}
%#############################################################

Here we discuss a fully non-perturbative renormalization prescription that will remove all divergences inherited in the matrix elements
of the quasi-PDFs, as well as the mixing, as indicated in the perturbative analysis of Ref.~\cite{Constantinou:2017sej}. In a nutshell, the proposed renormalization program:
\begin{itemize}
\item[\bf 1] removes the linear divergence that resums into a multiplicative exponential factor,
$e^{-\delta m |z|/a + c |z|}$. The coefficient $\delta m$ represents the strength of the divergence and is
expected to be operator independent, as it is related only to the Wilson line. $c$ is an arbitrary scale
\cite{Sommer:2015hea} that can be fixed by such a renormalization prescription;
\item[\bf 2] takes away the logarithmic divergence with respect to the regulator, $\log(a\bar\mu_0)$, where $\bar\mu_0$ is
the RI$'$ renormalization scale;
\item[\bf 3] applies the necessary finite renormalization related to the lattice regularization;
\item[\bf 4] eliminates the mixing that appears in the unpolarized operator, as the bare matrix element is
a linear combination of the unpolarized quasi-PDF and the twist-3 scalar operator. The two may be disentangled by
the construction of a $2{\times}2$ mixing matrix.
\end{itemize}

\medskip
We adopt a renormalization scheme which is applicable non-perturbatively, that is, the RI$'$~scheme \cite{Martinelli:1994ty}.
We compute vertex functions of the operators under study, between external quark states, with the setup being in momentum
space, and the operator defined as:
\be
{\cal O}_\Gamma = \overline \psi(x) \, \Gamma \,\mathcal{P}\, e^{i\,g\,\int_{0}^{z} A(\zeta) d\zeta}\, \psi(x+z\hat{\mu})\,,
\ee
where $\Gamma=\gamma_\mu,\,\gamma_\mu\cdot\gamma_5,\,\sigma_{\mu\nu}$ ($\nu\neq\mu$). The path ordering of the exponential
appearing in the above expression becomes, on the lattice, a series of path ordered gauge links. The renormalization functions
($Z$-factors) depend on the length of the Wilson line and, thus, we perform a separate calculation for each value of $z$. Typically, $z$
goes up to half of the spatial extent of the lattice.

\medskip
The renormalization prescription is along the lines of the program developed for local operators and the construction of the vertex functions is
described in Ref.~\cite{Alexandrou:2010me}. The difference between the renormalization of the local operators and
the Wilson-line operators is the linear divergence that appears in the latter case. However, there is no need to separate this divergence
from the multiplicative renormalization and, therefore, the technique described below may successfully extract both contributions
at once.

%%%%%%%%%%%%%%%%%%%%%%%%%%%%%%%%%%%%%%%%%%%%%
\vspace*{0.55cm}
%\newpage
\centerline{\bf\underline{\textit{Helicity and transversity quasi-PDFs}}}
\vskip 0.25cm
%%%%%%%%%%%%%%%%%%%%%%%%%%%%%%%%%%%%%%%%%%%%%

We first provide the methodology for a general operator with a Wilson line in the absence of any mixing. This is applicable for the
helicity and transversity quasi-PDFs, provided that their Dirac structure is chosen along the Wilson line. The renormalization functions of the
Wilson-line operators, $Z_{\cal O}$,
are extracted by imposing the following conditions:
\be
Z_q^{-1}\,Z_{\cal O}(z)\,\frac{1}{12} {\rm Tr} \left[{\cal V}(p,z) \left({\cal V}^{\rm Born}(p,z)\right)^{-1}\right] \Bigr|_{p^2=\bar\mu_0^2} = 1\, ,
\label{renorm}
\ee
where $Z_q$ is the renormalization function of the quark field obtained via
\be
Z_q = \frac{1}{12} {\rm Tr} \left[(S(p))^{-1}\, S^{\rm Born}(p)\right] \Bigr|_{p^2=\bar\mu_0^2}  \,.
\label{Zq_cond}
\ee
The trace is taken over spin and color indices, and the momentum $p$ entering the vertex function is set to the RI$'$ renormalization scale
$\bar\mu_0$. In Eq.~(\ref{renorm}) ${\cal V}(p,z)$ is the amputated vertex function of the operator and ${\cal V}^{{\rm Born}}$ is its
tree-level value, i.e. ${\cal V}^{{\rm Born}}(p,z) {=} i\gamma_3\gamma_5 \,e^{i p z}$ for the helicity operator. Also, $S(p)$ is the fermion
propagator and $S^{{\rm Born}}(p)$ is its tree-level. The RI$'$ scale $\bar\mu_0$ is chosen such that its $z$-component is
the same as the momentum of the nucleon. Such a choice serves as a suppression of discretization effects, as different classes of spatial
components have different discretization effects, and scales of the form $(n_t,3,3,3)$ have small discretization effects
\cite{Alexandrou:2015sea}. We test both diagonal (democratic) and parallel momenta to the Wilson line (in the spatial direction), that is
$a\vec{\bar{\mu}}_0{=} \frac{2\pi}{L}\,(P_3,P_3,P_3)$ and $a\vec{\bar{\mu}}_0{=} \frac{2\pi}{L}\,(0,0,P_3)$, respectively. We will refer
to these choices as ``diagonal'' and ``parallel''. The latter are expected to have larger lattice artifacts, as the ratio
\be
\label{Phat}
\hat{P}{\equiv}\frac{\sum_\rho \bar\mu_{0_\rho}^4}{\left(\sum_\rho \bar\mu_{0_\rho}^2\right)^2}
\ee
 is higher than for diagonal momenta. Using renormalization scales leading to a small value for such a ratio has been successful for the local
fermion operators~\cite{Constantinou:2010gr,Alexandrou:2010me}.

\medskip
The vertex functions ${\cal V}(p)$ contain the same linear divergence as the nucleon matrix elements.
This is crucial as it allows the extraction of the exponential together with the multiplicative $Z$-factor,
through the renormalization condition of Eq.~(\ref{renorm}). Regarding the renormalizability of quasi-PDFs,
it was proven to be multiplicative to all orders in Refs.~\cite{Ishikawa:2016znu,Ishikawa:2017faj}.
The authors consider quasi-PDFs defined in coordinate space and prove the multiplicative renormalizability,
and the same holds for the definition in Eq.~(\ref{MElements}). This is due to the fact that the former require,
in addition, a consistent subtraction scheme to remove all terms divergent in the limit $\xi_z\rightarrow 0$.
Moreover, the renormalization of bilocal composite operators is also studied in Ref.~\cite{Ishikawa:2016znu,Ishikawa:2017faj}.
Based on the above, $Z_{\cal O}$ can be factorized as
\be
Z_{\cal O}(z) = \overline{\cal Z}_{\cal O} \,\, e^{+\delta m |z|/a - c |z|}\,,
\ee
where $\overline{\cal Z}_{\cal O}$ is the multiplicative $Z$-factor of the operator and $\delta m$ is the strength of the linear divergence.
The exponential includes a term with an arbitrary scale $c$ that could be of the form $c|z|$~\cite{Sommer:2015hea}. Note that the exponential comes with a
different sign compared to the nucleon matrix element, as $Z_{\cal O}$ is related to the inverse
of the vertex function,
\be
Z_{\cal O} = \frac{Z_q}{\frac{1}{12} {\rm Tr} \left[{\cal V}(p) \left({\cal V}^{\rm Born}(p)\right)^{-1}\right] \Bigr|_{p=\bar\mu} }\,.
\ee
Such a construction of the $Z$-factor justifies the reason why the elimination of the power divergence in the nucleon matrix elements is successful
by multiplying with $Z_{\cal O}$, provided that it has been calculated on the same ensemble.
Consequently, the above renormalization condition handles all the divergences which are present in the matrix element under
consideration. Note that in the absence of a Wilson line ($z{=}0$), the $Z$-factors reduce to the ones for local currents,
which are free of any power divergence. For example, for the helicity operator, $Z_A(z{=}0)$ is the standard $Z$-factor
of the axial current.

\medskip
We would like to point out that knowledge of the coefficient $\delta m$ provides insight on the strength of the power divergence.
One can pursue this direction via the static potential~\cite{Ishikawa:2016znu} or the technique proposed in Ref.~\cite{Constantinou:2017sej}.
If the linear divergence is extracted, one may apply the matching of Ref.~\cite{Chen:2016fxx}, which includes the coefficient $\delta m$.
However, there is still a necessity to compute the multiplicative $Z$-factor to cure any logarithmic divergences, apply finite
renormalization, as well as fix the arbitrary scale $c$.


%%%%%%%%%%%%%%%%%%%%%%%%%%%%%%%%%%%%%%%%%%%%%
\vspace*{0.75cm}
\centerline{\bf\underline{\textit{Unpolarized quasi-PDF}}}
\vskip 0.25cm
%%%%%%%%%%%%%%%%%%%%%%%%%%%%%%%%%%%%%%%%%%%%%

The case of the unpolarized quasi-PDF requires special treatment, if the Dirac structure is in the same direction
as the Wilson line. As demonstrated in Ref.~\cite{Constantinou:2017sej}, for such a choice there is a mixing with
the twist-3 scalar operator\footnote{For twisted mass fermions the mixing is between the vector and pseudoscalar
currents.}, that must be eliminated. This mixing appears in some lattice regularizations (in particular, Wilson and twisted mass fermions) due to the breaking of chiral symmetry, and is found to be finite.
We establish notation by using the subscripts $S$ and $V$ for the scalar and vector (unpolarized) operators, respectively.
The corresponding operators are:
\bea
{\cal O}_S \hspace*{-0.15cm}&=&\hspace*{-0.15cm} \overline \psi(x) \, \hat{1} \,\mathcal{P}\, e^{i\,g\,\int_{0}^{z} A(\zeta) d\zeta}\, \psi(x+z\hat{\mu})\,, \\[2ex]
{\cal O}_V \hspace*{-0.15cm}&=&\hspace*{-0.15cm} \overline \psi(x) \, \gamma_\mu \,\mathcal{P}\, e^{i\,g\,\int_{0}^{z} A(\zeta) d\zeta}\, \psi(x+z\hat{\mu})\,,
\eea
and we represent their nucleon matrix elements as $h_S(P_3,z)$ and $h_V(P_3,z)$.
To disentangle the two contributions from their bare matrix elements, one must compute the multiplicative renormalization
and mixing coefficients from the following $2{\times}2$ matrix:
\begin{equation}
  \binom{{\cal O}^R_V(P_3,z)}{{\cal O}^R_S(P_3,z)} = \hat{Z}(z)\cdot
  \binom{{\cal O}_V(P_3,z)}{{\cal O}_S(P_3,z)}\,,
\end{equation}
where $\hat{Z}$ is the matrix of the mixing between the scalar and vector operators, that is
\begin{equation}
 \hat{Z}(z) = \begin{pmatrix} Z_{VV}(z) & Z_{VS}(z) \\ Z_{SV}(z) & Z_{SS}(z) \end{pmatrix} \,.
\end{equation}
%
According to the above mixing, the renormalized unpolarized quasi-PDF, $h^R_V(P_3,z)$, is related to the bare scalar and unpolarized via:
%
\begin{equation}
h^R_V(P_3,z) = Z_{VV}(z) \,\, h_V(P_3,z) + Z_{VS}(z) \,\, h_S(P_3,z) \,,
\label{h_R}
\end{equation}
%
where $Z_{VV}$ and $Z_{VS}$ are computed in a certain scheme. In the present work
we employ the RI$'$ scheme and then convert to the $\MSb$ scheme, at an energy scale $\bar\mu$.
The $Z_{ii}$ factors can be computed following a prescription similar to Eq.~(\ref{renorm}), that is:
\be
\label{renormMIX}
Z_q^{-1}\,\hat{Z}(z)\, {\cal \hat{V}}(p,z)\Bigr|_{p=\bar\mu} = \hat{1}\,,
\ee
%
where the elements of the vertex function matrix ${\cal \hat{V}}$ are given by the trace
\be
\left({\cal \hat{V}}(z) \right)_{ij}=  \frac{1}{12} {\rm Tr} \left[{\cal V}_i(p,z) \left({\cal V}_j^{\rm Born}(p,z)\right)^{-1}\right]\,,\quad i,j=S,V\,.
\ee
In the above equation ${\cal V}_i^{\rm Born}$ is the tree-level expression of the operator ${\cal O}_i$.
Thus, all matrix elements of $\hat{Z}$ can be extracted by a set of linear equations, which can be written in
the following matrix form:
\be
Z_q^{-1}\,\begin{pmatrix} Z_{VV}(z) & Z_{VS}(z) \\[2ex] Z_{SV}(z) & Z_{SS}(z) \end{pmatrix}  \cdot
\begin{pmatrix}  \left({\cal \hat{V}}(z) \right)_{VV} & \left({\cal \hat{V}}(z) \right)_{SV} \\[2ex]
 \left({\cal \hat{V}}(z) \right)_{VS} &\left({\cal \hat{V}}(z) \right)_{SS} \end{pmatrix} =
\begin{pmatrix} 1& 0 \\[2ex] 0 & 1\end{pmatrix} \,.
\ee
As can be seen from Eq.~(\ref{h_R}), a major component in the renormalization of
the unpolarized quasi-PDF is knowledge of the scalar nucleon matrix element $h_S(P_3,z)$. To date, no lattice calculation
is available for $h_S(P_3,z)$, as a mixing was not anticipated prior the work of Ref.~\cite{Constantinou:2017sej}. Therefore,
a proper renormalization of the unpolarized quasi-PDF is still pending. However, the mixing appears to be greatly suppressed
in the presence of a clover term in the fermionic action. This is of high importance, as we are currently computing the quasi-PDFs on
an ensemble of twisted mass clover improved fermions at the physical pion mass~\cite{Abdel-Rehim:2015pwa,Abdel-Rehim:2015owa}.


%#############################################################
\vspace*{0.25cm}
\subsection{Conversion to the $\MSb$-scheme}
\vspace*{0.15cm}
%#############################################################

The fermion operators with a finite-length Wilson line are scheme and scale dependent. As a result, having obtained
the $Z$-factors in the RI$'$ scheme as depicted in Eqs.~(\ref{renorm}) and (\ref{renormMIX})
at the RI$'$ scale $\bar\mu_0$, we must convert them to the $\MSb$ scheme at a scale $\bar\mu$. This conversion factor has
been computed in one-loop continuum perturbation theory in Ref.~\cite{Constantinou:2017sej}. In addition, comparison
with phenomenological estimates is done typically at $\bar\mu{=}2$ GeV. This defines the value of the $\MSb$ renormalization
scale chosen in the expressions for the conversion factors. Alternatively, one can choose $\bar\mu{=}\bar\mu_0$, and then
evolve the scale to 2 GeV, which requires the anomalous dimension of the operator:
\be
\gamma_{\cal O} =  -3 \frac{g^2\,C_f}{16\pi^2} \,.
\ee
Note that to one-loop level, the anomalous dimension does not dependent on the Dirac structure of the operator
and is the same in the RI$'$ and $\MSb$ schemes. The evolution to the $\MSb$ renormalization scale
$\bar\mu{=}2$ GeV can be performed using the intermediate Renormalization Group Invariant scheme (RGI), as
employed in Refs.~\cite{Constantinou:2014fka,Alexandrou:2015sea}. In the framework of this paper we tested both
methods and their difference was found to be negligible.

\begin{figure}[h]
\centering
\includegraphics[scale=0.265,angle=-90]{CV_CA_RIMS_Landau_pL3_Re.eps}\,\,\,
\includegraphics[scale=0.265,angle=-90]{CV_CA_RIMS_Landau_pL3_Im.eps}
\vspace*{-0.3cm}
\begin{minipage}{15cm}
\hspace*{3cm}
\caption{\small{One-loop conversion factor from the RI$'$ to the $\MSb$ scheme for $\bar\mu{=}\bar\mu_0$. In
the left (right) plot we show the real (imaginary) part of the conversion factor as a function
of the length of the Wilson line in lattice units. Two choices of the RI$'$ scale have been
employed: $a\bar\mu_0{=}\frac{2\pi}{32}\, (\frac{7}{2}{+}\frac{1}{4}, 3,3,3)$ (blue circles)  and
$a\bar\mu_0{=}\frac{2\pi}{32}\, (\frac{4}{2}{+}\frac{1}{4}, 0,0,3)$  (red squares).
We use the abbreviation (7,3,3,3) and (4,0,0,3) in the legends, respectively. }}
\label{fig:C}
\end{minipage}
\end{figure}

\medskip
As discussed in the previous section, we employ two types of the RI$'$ renormalization scale regarding the spatial direction:
one the same as in the nucleon matrix elements, $a\vec{\bar{\mu}}_0{=}\frac{2\pi}{L}\,(0,0,P_3)$, (parallel to the Wilson line direction),
and one which is diagonal and each direction has a value of $P_3$, that is $a\vec{\bar{\mu}}_0{=}\frac{2\pi}{L}\,(P_3,P_3,P_3)$.
The conversion factor depends on both the RI$'$ and $\MSb$ renormalization scales. While $\bar\mu$ is typically fixed to 2 GeV,
$\bar\mu_0$ can change affecting the numerical values of the conversion factor. Such a case is illustrated in Fig.~\ref{fig:C}
for the unpolarized and helicity operators, which share the same conversion factor\footnote{The one-loop calculation does not
depend on the prescription which one adopts for extending $\gamma_5$ to $D$ dimensions (see, e.g.,
Refs.~\cite{Buras:1989xd,Patel:1992vu,Larin:1993tp,Larin:1993tq,Skouroupathis:2008mf,Constantinou:2013pba}
for a discussion of four relevant prescriptions and some conversion factors among them).}.
We focus on the renormalization of the matrix elements with the nucleon boosted by $a P_3{=}\frac{6\pi}{L}$, and we apply the same
to the conversion factor, for the two cases of the renormalization scale. For the specific value of $P_3$, we choose the temporal
direction of $\bar\mu_0$ such that the ratio $\hat{P}$ defined above, is as small as possible in order to suppress
lattice artifacts. Nevertheless, for the ``parallel'' case the minimum value for the ratio is $\hat{P}{=}0.54$ which is already very high.
The chosen values for $a\bar\mu_0$ are: $\frac{2\pi}{32}\, (\frac{7}{2}{+}\frac{1}{4}, 3,3,3)$ and
$\frac{2\pi}{32}\, (\frac{4}{2}{+}\frac{1}{4}, 0,0,3)$ for the ``diagonal'' and ``parallel'' case, respectively.

\medskip
The real and imaginary parts of the conversion factor are plotted as a function of the length of the Wilson
line, $z$. We stress that the dependence of the conversion factor on the renormalization scales and the length of the
Wilson line is highly non-trivial and is expressed in terms of integrals of modified Bessel functions. Consequently, the data points
shown in Fig.~\ref{fig:C} have been computed numerically for the specific scales, at each value of $z$ separately. We observe
that the real part of the conversion is an order of magnitude larger than the imaginary part. The real part consistently increases
for increasing values of $z$, while the imaginary part almost immediately stabilizes when $z$ becomes non-zero. Also, the conversion
factor at $z{=}0$ is equal to unity, as it corresponds to the local vector and axial currents. Note that the case $z{=}0$ is not extracted
from the calculation of Ref.~\cite{Constantinou:2017sej} as it is strictly for $z{\neq}0$\,: the appearance of contract terms beyond
tree level renders the limit $z\to 0$ nonanalytic. On the contrary, the non-perturbative prescription of the previous section
can be applied for any value of $z$, as the calculation is performed on each $z$ separately. The values of Fig.~\ref{fig:C} will be used
in the following section to bring the RI$'$ non-perturbative $Z$-factors to the $\MSb$ scheme. To reliably extract
the $Z$-factors we have extend the calculation including several values of $\bar\mu_0$ as explained in Section~\ref{Sec3}. The conversion
factor was found to have the same qualitative behavior for all values of $\bar\mu_0$.
